The four methods used in this report provide complementary information on the same empirical question. The linear time-series analysis highlights the strong hourly and weekly structure of the market. The extreme-value analysis shows that rare and economically meaningful stress episodes remain an essential feature of the 2025 price distribution. The copula analysis reveals that dependence between France and Germany is not fully summarized by linear correlation and becomes especially informative in the tails. Finally, the Hawkes analysis shows that extreme events are temporally clustered and partially transmitted across markets.

Taken together, these findings suggest that market stress in the French-German electricity system has both a cross-sectional and a temporal dimension. Dependence matters because prices can become jointly extreme, and self-excitation matters because extreme events tend to arrive in bursts rather than in isolation. This joint perspective is one of the main contributions of the report.
